\subsection{Intersection-closed ample classes and signed bouquets}
\label{sec:signed-bouquets}

Intersection closure singles out a classical family with explicit allocation
maps. It is larger than the convex-geometry ideals already appearing in the
atlas. This section consolidates that distinction and the relevant recent
results; the longer operator proofs remain in their original notes.

\begin{definition}[Bouquets and their signed version]
\label{def:atlas-signed-bouquets}
A \emph{bouquet of convex geometries} on a finite set $E$ is a family
$C\subseteq2^E$ such that $\varnothing\in C$, $A\cap B\in C$ for all
$A,B\in C$, and
\[
 A,B\in C,\quad A\subsetneq B
 \quad\Longrightarrow\quad
 A\cup\{e\}\in C\text{ for some }e\in B\setminus A.
\]
The last property is local one-element extendability. These are also called
\emph{conditional antimatroids}. A bouquet containing $E$ is a convex
geometry. Define
\[
 \mathsf{Bouq}^{\pm}
 :=\{C\in\Amp:\text{some }C\mathbin\triangle r
                  \text{ is intersection-closed}\},\qquad
 C\mathbin\triangle r:=\{A\mathbin\triangle r:A\in C\}.
\]
The superscript records coordinate flips, not ambient cube complementation.
\end{definition}

\begin{proposition}[Characterization and minor closure]
\label{prop:atlas-bouquet-characterization}
For a nonempty intersection-closed family containing $\varnothing$,
ampleness is equivalent to local one-element extendability. Consequently
$\mathsf{Bouq}^{\pm}$ is exactly the signed-coordinate class of bouquets,
with constant coordinates suppressed as usual. It is p-closed and c-closed,
and hence pc-closed. For an ample $C$, equivalently
\[
 C\in\mathsf{Bouq}^{\pm}
 \quad\Longleftrightarrow\quad
 \exists r\in C\quad
 \operatorname{maj}(u,v,r)\in C\quad(u,v\in C),
\]
where $\operatorname{maj}$ denotes coordinatewise majority.
\end{proposition}
\begin{proof}
An ample class is isometric. A Hamming geodesic between comparable members
adds their differing elements one at a time, giving local extendability.
The converse is the classical ampleness theorem for conditional antimatroids
\cite[\S3.2.4]{CCMW}; the equivalent bouquet axioms are
\cite[\S2.5]{CCKGeometry}.

For any nonempty intersection-closed family $H$, its total intersection
$m=\bigcap H$ is a member. Flipping the coordinates of $m$, which are
constant one on $H$, gives an intersection-closed family containing zero.
Thus an initially arbitrary witnessing flip can be changed to an actual
word $r\in C$. In the orientation with $r=0$, the displayed majority
condition is precisely closure under intersection, proving the equivalence.

Projection commutes with intersection. A conditioning is also
intersection-closed: the meet of two words agreeing with the fixed pattern
still agrees with it. Both operations preserve ampleness. Coordinate flips
restrict to the remaining coordinates; newly constant one coordinates can
be normalized by the total-intersection argument. This proves both closure
claims, including conditionings that discard the original root.
\end{proof}

\begin{example}[Why the full ground set is not required]
\label{ex:atlas-bouquet-star}
On three coordinates, $\{\varnothing,\{1\},\{2\},\{3\}\}$ is a bouquet
and a daisy cube, but is not a convex geometry on that ground set. It is
nevertheless an ideal of the full Boolean convex geometry. Thus omission
of the full set alone does not distinguish bouquets from convex-geometry
ideals; \cref{prop:atlas-bouquet-position} gives a genuine separation.
\end{example}

\paragraph{Recognition by forbidden partial patterns.}
A \emph{minimally missing pattern} $(S,a)$ has
$a\notin C|_S$, while every proper restriction of $a$ occurs. For a proposed
flip $r$, intersection closure is equivalent to
\begin{equation}\label{eq:atlas-bouquet-horn}
 \bigl|\{i\in S:r_i=a_i\}\bigr|\le1
 \qquad\text{for every minimally missing }(S,a).
\end{equation}
Indeed, agreement at two positions lets two witnesses for the proper
restrictions have a majority with $r$ realizing the missing pattern.
Conversely, if $w=\operatorname{maj}(u,v,r)$ is absent, choose a minimally
missing restriction of $w$. The nonempty sets of positions where $u$ and
$v$ disagree with it are disjoint, and $r$ agrees with it on both sets,
contradicting \eqref{eq:atlas-bouquet-horn}.
The pairwise prohibitions are a 2-SAT formula in the bits of $r$. This is
the classical \emph{renamable-Horn} condition, not a new recognition
theorem; see \cite[\texttt{lem:horn-test}]{BouquetProjection}.
The 2-SAT size is measured in the supplied minimally missing patterns,
not in an unspecified implicit description of $C$.

\paragraph{Explicit maps on both complementary classes.}
Work in an intersection-closed orientation $H$ on the full original
ground set $E$. Put
\[
 \begin{aligned}
 \rho(v)&=\{i\in v:v\setminus\{i\}\in H\},&
 \downarrow H&=\{a:a\subseteq v\text{ for some }v\in H\},\\
 c(a)&=\bigcap\{v\in H:a\subseteq v\}
       &&(a\in\downarrow H).
 \end{aligned}
\]
The classical representation map is $v\mapsto\rho(v)$; its whole downward
cube is present by intersection closure. Deleting words in decreasing
cardinality gives the classical corner peeling
\cite[Proposition~4.9]{CCMW}. The closure fibers satisfy
\begin{equation}\label{eq:atlas-bouquet-fibers}
 c^{-1}(v)=[\rho(v),v]
 :=\{a:\rho(v)\subseteq a\subseteq v\}.
\end{equation}
If $c(a)=v$, omitting a removable element contradicts minimality of $c(a)$.
Conversely, if $\rho(v)\subseteq a\subseteq v$ but $c(a)\subsetneq v$,
a geodesic from $c(a)$ to $v$ ends at a lower neighbor $v\setminus\{i\}$,
forcing $i\in\rho(v)\setminus c(a)$, a contradiction.

The recent complementary construction in
\cite[\texttt{thm:closure-completion}]{BouquetCompletion} assigns, on
$K=2^E\setminus H$,
\begin{equation}\label{eq:atlas-bouquet-complement}
 r_K(a)=
 \begin{cases}
 E\setminus a,&a\notin\downarrow H,\\
 (E\setminus c(a))\cup(a\setminus\rho(c(a))),
       &a\in\downarrow H\setminus H.
 \end{cases}
\end{equation}
It is an ordinary representation map made of actual whole cubes. More
strongly, it is the remaining-neighbor map of a corner peeling: process
closure fibers by increasing $|v|$, process each
$[\rho(v),v]\setminus\{v\}$ by decreasing $|a|$, and finally process
$2^E\setminus\downarrow H$ by increasing $|a|$.
The cited proof verifies every selected cube and every omitted neighboring
direction. It retains original constants and allows the empty complement.
This is an author-proved recent result, with external priority unresolved;
it is not attributed to the classical primal-map theorem.

\begin{proposition}[Position in the atlas]
\label{prop:atlas-bouquet-position}
With the preceding complementary-completion input,
\[
 \mathsf{Icg}\subsetneq\mathsf{Bouq}^{\pm}
 \subseteq\mathcal B\subsetneq\mathsf{Ext},\qquad
 \mathsf{Bouq}^{\pm}\not\subseteq\RAmp.
\]
In particular, signed bouquets are corner-peelable and representable.
\end{proposition}
\begin{proof}
An ideal of a convex geometry is closed under intersection, since the
intersection is a closed set below either member of the ideal. It is ample
by \cite[\S2.5]{CCKGeometry}, so signed ideals lie in
$\mathsf{Bouq}^{\pm}$. The two peeling constructions give
$\mathsf{Bouq}^{\pm}\subseteq\mathcal B$.
The strict inclusion $\mathcal B\subsetneq\mathsf{Ext}$ is already proved
in \cref{prop:maxmin-neighbors}.

The 26-set bouquet of \cite{NonrealizableBouquet} has no affine-convex
realization in any finite dimension. Every signed convex-geometry ideal
is realizable by \cite[Theorem~26]{CCKGeometry}, and realizability is
invariant under coordinate flips. This proves both the first strictness
and the displayed noncontainment. That same bouquet has a certified
maximum parent; it is not a non-$\MaxMin$ example.
\end{proof}

\paragraph{Attaching a cube without assuming a daisy intersection.}
The complementary construction extends beyond a single bouquet.
\begin{proposition}[A clamped cube on a signed bouquet]
\label{prop:atlas-bouquet-cube-attachment}
Let $B\in\mathsf{Bouq}^{\pm}$ be nonempty on the original frame $E$, and
let $C=Q(z,S)$ be a positioned cube. Suppose its coordinate clamp
\[
 c_C(b)|_S=b|_S,\qquad c_C(b)|_{E\setminus S}=z|_{E\setminus S}
\]
satisfies $c_C(B)\subseteq B$. Then $B\cup C\in\mathsf{Ext}$.
Every supplied ordinary representation map on $B$ extends to $B\cup C$
without changing its old rooted-cube columns; an ordinary map on
$(B\cup C)^c$ can be constructed simultaneously. Cyclic supplied maps
are allowed. No chordality or daisy-intersection hypothesis is required.
\end{proposition}
\begin{proof}
Flip coordinates to make $B$ intersection-closed. Clamps commute with
this change of coordinates. The compatible-completion theorem
\cite[\texttt{thm:horn-gluing}]{CubeComplementCompletion} constructs a full
actual-cube basis, retains the supplied primal basis, and supplies the
entire complementary basis. Translate its two maps back to the original
frame, retaining all constants.
\end{proof}
The union need not itself be a signed bouquet. In particular this is an
attachment theorem, not a closure law for $\mathsf{Bouq}^{\pm}$, nor an
iteration rule for arbitrary unions of differently oriented bouquets.

\paragraph{Operator characterization and remaining question.}
For $g_0=(1,0)^t$ and $g_1=(1,-1)^t/\sqrt2$, let $P_C$ be the orthogonal
projection onto the span of the tensors $\bigotimes_{i\in E}g_{v_i}$,
$v\in C$. CP's author theorem
\cite[\texttt{thm:classification}]{BouquetProjection} says, for isometric
$C$, that $P_C$ is entrywise nonnegative in some real orthonormal product
basis exactly when a flip makes $C$ intersection-closed. A product basis
chooses an orthonormal basis separately at each coordinate. Entrywise
nonnegativity is stronger than positive semidefiniteness. This is an
algebraic characterization of the same family, not evidence that all ample
classes admit such a basis. Its full spectral proof is not independently
certified by this atlas consolidation.

The question whether every signed bouquet has a finite maximum parent
remains open in the campaign. The ideal subclass has such parents, and the
nonrealizable 26-set example does too. Neither the explicit two-shore maps
nor the operator characterization supplies parents for all bouquets.
No comparison with all of $\Rper$ or all of $\mathsf{CAmp}$ is asserted.
